\begin{center}
\textsc{Gauss} an \textsc{Bessel}.\quad Göttingen, 1.~April 1827.
\end{center}

\par\noindent\hspace*{\parindent}.\ .\ .\ .\ .\ .\quad Ich möchte in diesem Frühjahr die Amplitude des Bogens zwischen
Göttingen und Altona mit dem Zenitsektor messen und bin selbst gespannt
auf das Ergebnis. Eine Reihe von Beobachtungen, die Anfang
1824 an den Meridiankreisen vorgenommen worden waren, ergab \(2^{0}\,1^{\prime}\,58^{\prime\prime}\). Die Polhöhen, meine nach den
besten Beobachtungen angesetzt und die von \textsc{Schumacher} nach seiner eigenen Angabe, ergeben
eher 1 oder \(1\frac{1}{2}\) Bogensekunden weniger. Die geodätischen Messungen hingegen
ergeben unter der Annahme der gleichförmigen Gestalt der Erde und mit \textsc{Walbecks} Dimen-
\par\clearpage
\noindent{\small\textbf{ZUR HANNOVERSCHEN TRIANGULATION.}}\hfill363\par\medskip
\noindent sionen 4 oder 5 Bogensekunden mehr. Fast genau ebenso viel wie dieses
letztere Ergebnis liefern die ganz gleichzeitigen Zenitdistanzen von Zenitsternen,
die ich hier und \textsc{Nehus} in Altona am Meridiankreis beobachtet haben, wobei
der Nullpunkt mit Kollimatoren bestimmt wurde. \textsc{Nehus} ist jetzt hier, um mir bei
den Sektorbeobachtungen zu helfen. Er hat einen hängenden \textsc{Repsold}schen Kolli\-%
mator mitgebracht, doch bisher halte ich die Methode, den Nadirpunkt durch
Quecksilberreflexion zu bestimmen, für bedeutend genauer. Ich denke aber, dass sich jene Methode
noch verbessern lässt.\quad\mbox{.\ .\ .\ .\ .\ .}
